\documentclass{report}
\usepackage{amsmath,amssymb}
\setlength{\parindent}{0mm}
\setlength{\parskip}{1em}
\begin{document}
\begin{center}
\rule{6in}{1pt} \
{\large Lin Lin \\
{\bf Accelerating the convergence of the SCF Iteration through charge mixing and preconditioning}}

Lawrence Berkeley National Laboratory \\ 1 Cyclotron Road Mail Stop 50F-1650 \\ Berkeley \\ CA 94720
\\
{\tt linlin@lbl.gov}\\
Chao Yang\end{center}

Kohn-Sham density functional theory (KSDFT) is the most widely used
electronic structure theory for condensed matter systems. Although KSDFT
is often stated as a nonlinear eigenvalue problem, an alternative
formulation of the problem, which is more convenient for understanding
the convergence of numerical algorithms for solving this type of problem,
is based on a nonlinear map known as the Kohn-Sham map. The solution to
the KSDFT problem is a fixed point of this nonlinear map. The simplest
way to solve the KSDFT problem is to apply a fixed point iteration to the
nonlinear equation defined by the Kohn-Sham map. This is commonly known
as the self-consistent field (SCF) iteration in the condensed matter
physics and chemistry communities. However, this simple approach often
fails to converge. The difficulty of reaching convergence can be seen
from the analysis of the Jacobian matrix of the Kohn-Sham map, which we
will present in this talk. The Jacobian matrix is directly related to the
dielectric matrix or the linear response operator in the condense matter
community. We will show the different behaviors of insulating and
metallic systems in terms of the spectral property of the Jacobian
matrix. We discuss how to use these properties to approximate the
Jacobian matrix and to develop effective preconditioners to accelerate
the convergence of the SCF iteration. Approximations to the Jacobian can
also be constructed directly using the Broyden-like technique. In
computational condense matter physics, this is often known as the charge
mixing scheme. We will present and compare different mixing schemes and
point out some open questions regarding the effects of mixing and
preconditioning on the convergence of the SCF iteration.


\end{document}
